\begin{declarations}

\begin{acknowledgements}
Withheld for double-anonymous review, in accordance with the journal's policy.
\end{acknowledgements}

\begin{contributions}
% Anonymised: iSys double-anonymous review forbids identifying the authors here.
The author(s) of this submission made the following contributions, stated by role rather than by
name to preserve anonymity for review, in accordance with the CRediT taxonomy:
Conceptualization, Methodology, Software, Validation, Formal analysis, Investigation, Data curation,
Writing -- original draft, Writing -- review and editing, Visualization, Supervision, Project
administration. A named allocation of these roles will be supplied in the camera-ready version, as
required by the journal after acceptance.
\end{contributions}

\begin{interests}
The authors declare that they have no competing interests. No funding source had any role in the
design of the study, the collection, analysis or interpretation of data, or the decision to submit.
\end{interests}

\begin{materials}
The replication package --- including the sampling frame with its snapshot hash and seed, the
screening and attrition tables, the static-indicator and mutation harnesses with pinned tool versions,
the raw per-mutant statuses, the analysis scripts and the append-only run log --- is prepared for
public release in a trusted repository. Because iSys uses double-anonymous review, the package is
made available to reviewers through an anonymised mirror, and the canonical public repository URL is
supplied in the camera-ready version. All analysed repositories are public and were used under their
respective permissive licences, which are recorded per repository in the screening table; no
repository content is redistributed in the package. Third-party papers referenced in the literature
review are cited, not redistributed.
\end{materials}

\begin{funding}
This research received no specific grant from any funding agency in the public, commercial or
not-for-profit sectors. This statement will be updated if that changes before publication.
\end{funding}

\begin{furtherinformation}
\paragraph{Ethical considerations.}
This study analyses only publicly available software artefacts and involves no human participants, so
no ethics committee approval was required or sought. Mutation analysis modifies code only in isolated
working copies; no public repository was altered. No code was injected into any third-party
repository, and the seeded-drift validation described in the protocol was \emph{not} carried out in
this study (Section~\ref{sec:threats}); had it been, it would have been confined to separate copies so
that no measured artefact was modified. Repositories lacking a detectable permissive licence were
excluded before any modification.

\paragraph{Declaration of AI-Assisted Use in Writing.}
In accordance with the journal's policy on AI/LLM use, and in the interest of full transparency, we
declare the following. Generative AI tools were used in this work as follows: (i) as a coding and
drafting assistant for software engineering tasks in the replication package (writing and refactoring
harness scripts, documenting code, and generating tables and figures from result files); and (ii) for
drafting and editing prose in this manuscript, including restructuring sections and improving clarity
and language.

In every case the tools used were large language model assistants, and the following verification
applies. No AI tool is listed as an author, and none is credited with intellectual contribution. No
AI tool was used to generate research data, to invent measurements, or to fabricate citations. Every
quantitative value reported in this paper is produced by a command recorded in the replication
package's append-only run log, and every reference was verified against Crossref by an automated
audit before the manuscript was built; references that could not be verified were removed rather than
retained. The human author(s) reviewed, verified and take full responsibility for all content,
including all AI-assisted text.

We state one further, material use because it affects the interpretation of our results. AI tools
assisted in executing the measurement pipeline --- running screens, parsing tool output, and
assembling tables. They were not used as \emph{raters}. Where the protocol required human judgement
(the business-rule relevance of mutants, the mirror-step classification, and inter-rater reliability),
no AI labelling was substituted, and the absence of those reliability statistics is disclosed as a
limitation in Section~\ref{sec:threats} rather than filled with an automatically generated number.
Table~\ref{tab:ai-summary} summarises each use of generative AI and the verification applied.

\par\medskip
\noindent\begin{minipage}{\columnwidth}
\centering
\footnotesize
\captionof{table}{Summary of generative-AI use and the verification applied to each use.}
\label{tab:ai-summary}
\begin{tblr}{width=\columnwidth,colspec={X[0.72,l]X[1.68,l]},rowsep=2pt,colsep=4pt}
\toprule
\textbf{Purpose} & \textbf{Verification applied} \\
\midrule
Replication-package code and scripts & Executed; outputs checked against tool reports; run log records command, timestamp, commit and versions \\
Drafting and editing manuscript prose & Human author reviewed and approved all text; claims checked against result files \\
Tables and figures & Generated programmatically from result files, never typed by hand \\
Reference verification & Automated Crossref audit of every DOI; unverifiable entries deleted \\
Research data & \textbf{Not generated by AI.} All measurements come from executed tools; human-judgement items were not automated \\
\bottomrule
\end{tblr}
\end{minipage}
\par\medskip
\end{furtherinformation}

\end{declarations}
